Scaling processes couple strongly nonlinear kinetics to hard operational
constraints, and the models available for them are structurally wrong in
ways that matter. In this work the wrongness is measured rather than
assumed: a classical fitted rate law mispredicts independent modern
kinetic data by a factor of six at the median. We present a control
architecture that treats such mismatch as a first-class design object.
A hybrid soft sensor combines a physics core with a learned residual
whose weight is set by its own leave-one-out credibility, so the
correction fades back to the prior whenever the data has not earned it.
A vertex-robust economic real-time optimizer then consumes both the
correction and its uncertainty, the latter as a multiplicative
constraint back-off. On a synthetic brine-line plant grounded in the
primary-literature chemistry of Philippine liquid-dominated geothermal
fields, the sensor reduces estimation error by 40\% relative to physics
alone and 48\% relative to a purely data-driven model under sparse
daily labels. In closed loop it is the only tested policy that remains
essentially violation-free in two regimes: benign random mismatch
(+33\% power uplift, zero violation days) and an adverse stress set at
the measured sixfold kinetic gap (0.5 post-warm-up violation days
against 17.2 for the robust open loop). Two pre-registered failures,
of zero-shot residual transfer and of scalar modifier adaptation
against multiplicative mismatch, motivated the design and are reported
as results. Every kinetic, solubility and sensing component is
validated out-of-sample against an independent 2024 experimental
dataset.
